\section*{第 \ref{ch:g11:drift-to-spreading} 章——从大陆漂移到海底扩张}

\begin{solution}{exo:g11:drift-to-spreading:1}
\omterm{def:g11:drift-to-spreading:spreading}{海底扩张}是指在洋中脊的脊顶处形成新洋底的
过程：熔融岩石上升并凝固，随后洋底
向洋脊两侧远离运动。\omterm{def:g11:drift-to-spreading:spreading-rate}{全扩张速率}是两板块
相互远离的速度；半扩张速率是洋底的一侧
远离轴线的速度。对称扩张时，\omterm{def:g11:drift-to-spreading:spreading-rate}{全扩张速率} $= 2
\times$ 半扩张速率。
\end{solution}

\begin{solution}{exo:g11:drift-to-spreading:2}
高于其\omterm{def:g11:drift-to-spreading:remanence}{居里温度}（约 \qty{580}{\celsius}）时，磁铁矿无法保持自身的
磁化强度；只有冷却到\omterm{def:g11:drift-to-spreading:remanence}{居里温度}以下，其颗粒才会
沿磁场方向磁化并保持该方向。再次加热至
\qty{700}{\celsius}，高于\omterm{def:g11:drift-to-spreading:remanence}{居里温度}，岩石失去其原有的
记录；再次冷却时，岩石记录下该新时期的磁场。
\end{solution}

\begin{solution}{exo:g11:drift-to-spreading:3}
声音向下传播再返回：$d = 1500 \times 5.2/2 = \qty{3900}{m}$。
\end{solution}

\begin{solution}{exo:g11:drift-to-spreading:4}
当前\omterm{def:g11:drift-to-spreading:reversal}{正极性}时段始于 78 万年前：至今已持续
\num{780000} 年。雅拉米洛事件持续了 $1.07 - 0.99 = \num{0.08}$
百万年，即 \num{80000} 年。统计年龄小于五
百万年的界线（0.78、0.99、1.07、1.77、1.95、2.14、2.15、2.581、3.04、3.11、
3.22、3.33、3.58、4.18、4.29、4.48、4.62、4.80、4.89、4.98）：共二十次倒转。
\end{solution}

\begin{solution}{exo:g11:drift-to-spreading:5}
$65\,000/40\,000 \approx 1.6$：洋脊总长约
绕地球一周半。
\end{solution}

\begin{solution}{exo:g11:drift-to-spreading:6}
$20/0.78 \approx \qty{25.6}{km}$ 每百万年，即半扩张速率
约为 \qty{2.6}{cm} 每年；\omterm{def:g11:drift-to-spreading:spreading-rate}{全扩张速率}约为 \qty{5.1}{cm} 每年。
\end{solution}

\begin{solution}{exo:g11:drift-to-spreading:7}
西侧：$45/3.58 \approx 12.6$~千米每百万年，\qty{1.26}{cm} 每年；
东侧：$44/3.58 \approx 12.3$~千米每百万年，\qty{1.23}{cm} 每年。
\omterm{def:g11:drift-to-spreading:spreading-rate}{全扩张速率}：$1.26 + 1.23 \approx \qty{2.5}{cm}$ 每年。两侧的半扩张速率
相差约 \qty{2}{\percent}：扩张几乎是对称的。
\end{solution}

\begin{solution}{exo:g11:drift-to-spreading:8}
\qty{1.9}{cm} 每年即 \qty{19}{km} 每百万年：$990/19 \approx 52$
百万年，略大于最老沉积物的 4900 万年
（约 \qty{6}{\percent}）。
\end{solution}

\begin{solution}{exo:g11:drift-to-spreading:9}
按每百万年 \qty{20}{km}：\qty{15.6}{km}、\qty{19.8}{km}、
\qty{21.4}{km}、\qty{35.4}{km} 和 \qty{39.0}{km}。自轴线起：从 0
至 \qty{15.6}{km} 为\omterm{def:g11:drift-to-spreading:reversal}{正极性}，至 \qty{19.8}{km} 为\omterm{def:g11:drift-to-spreading:reversal}{反极性}，至 \qty{21.4}{km}
为\omterm{def:g11:drift-to-spreading:reversal}{正极性}（雅拉米洛事件），至 \qty{35.4}{km} 为\omterm{def:g11:drift-to-spreading:reversal}{反极性}，至 \qty{39.0}{km}
为\omterm{def:g11:drift-to-spreading:reversal}{正极性}（奥杜威事件），其后为\omterm{def:g11:drift-to-spreading:reversal}{反极性}。
\end{solution}

\begin{solution}{exo:g11:drift-to-spreading:10}
$64 \times 10^6/4.6 \times 10^9 \approx \num{0.014}$，约
占地球年龄的 \qty{1.4}{\percent}。这支持了赫斯的论点：
洋底之所以年轻，是因为它不断被生成与消亡；
地球历史的大部分时期内形成的洋底都未能保存至今。
\end{solution}

\begin{solution}{exo:g11:drift-to-spreading:11}
任一时刻，洋底在轴线处以当时的极性形成，
随即一分为二，各一半以相同速度向两侧移动。同时形成的两条带
因此与轴线距离相同，分别位于两侧，
极性相同：两侧互为镜像，关于轴线对称。对称性
只有当一侧洋底远离的速度快于另一侧（不对称扩张），
或轴线跳至新
位置时，对称性才会被打破。
\end{solution}

\begin{solution}{exo:g11:drift-to-spreading:12}
(1) \qty{5}{cm} 即 $5 \times 10^{-5}$~\unit{km}：每年新增一条
面积为 $65\,000 \times 5 \times 10^{-5} = \qty{3.25}{km^2}$ 的带状区域。
(2) $2.9 \times 10^8/3.25 \approx 9 \times 10^7$ 年，约
9000 万年。(3) 这与最老洋底的年龄处于同一量级（
1.9 亿年），但更短：要使洋底保持如此年轻，
必须在某处消亡的洋底量与洋脊处生成的一样多。
\end{solution}

\begin{solution}{exo:g11:drift-to-spreading:13}
该事件持续了 8 万年。按每百万年 \qty{12.5}{km} 计，
其条带宽 $0.08 \times 12.5 = \qty{1}{km}$；若每百万年 \qty{75}{km}，
则宽 $0.08 \times 75 = \qty{6}{km}$。从海面观测，
海面上方 \qty{3}{km} 处，宽度远小于深度的条带
只产生微弱磁场，并与两侧相\omterm{def:g11:drift-to-spreading:reversal}{反极性}的宽条带磁场混在一起；
宽度为深度两倍的条带则清晰可见。
\end{solution}

\begin{solution}{exo:g11:drift-to-spreading:14}
站位 16：$176/11 \approx \qty{16}{m}$ 每百万年；站位 20：$72/67
\approx \qty{1.1}{m}$ 每百万年。沉积物堆积速率
随地点和时间差异很大，因此其厚度
不能用来确定洋底的年龄。洋底的年龄由直接覆盖在玄武岩上的
最老沉积物中的化石给出。
\end{solution}

\begin{solution}{exo:g11:drift-to-spreading:15}
(1) 西侧：$58/2.581 \approx \qty{22.5}{km}$ 每百万年，
\qty{2.25}{cm} 每年；东侧：$70/2.581 \approx \qty{27.1}{km}$ 每百万年，
\qty{2.71}{cm} 每年；\omterm{def:g11:drift-to-spreading:spreading-rate}{全扩张速率} $128/2.581 \approx \qty{49.6}{km}$ 每
百万年，约 \qty{5.0}{cm} 每年。(2) 西侧约 \qty{225}{km}，东侧约
\qty{271}{km}。(3) 反转发生时，两条边界在轴线处重合；
此后两者以\omterm{def:g11:drift-to-spreading:spreading-rate}{全扩张速率}相互远离，
无论运动在两侧如何分配。两者间距除以
其年龄，即为\omterm{def:g11:drift-to-spreading:spreading-rate}{全扩张速率}。
\end{solution}

\begin{solution}{pb:g11:drift-to-spreading:1}
\textbf{1.} 16 号站点最近（\qty{191}{km}，1100 万年），20 号站点
最远（\qty{1270}{km}，6700 万年）：离洋脊越远，
洋底越老。

\textbf{2.} 玄武岩一形成，沉积物便开始在其上沉积，因此最老的一层
与玄武岩年龄大致相同。它只可能
略年轻一些，绝不会更老。

\textbf{3.} 以千米每百万年为单位，16 号站点：$191/11 \approx 17.4$；
15 号站点：$380/24 \approx 15.8$；18 号站点：$506/26 \approx 19.5$；17 号站点：
$643/33 \approx 19.5$；14 号站点：$727/40 \approx 18.2$；19 号站点：$990/49
\approx 20.2$；20 号站点：$1270/67 \approx 19.0$。

\textbf{4.} $129.5/7 \approx \num{18.5}$~千米每百万年，约
\qty{1.85}{cm} 每年。

\textbf{5.} 是的：其比值（19.5）落在西侧各站点的比值范围之内。
两侧洋底扩张速率大致相同：扩张是
对称的。

\textbf{6.} $t = d/u$：年龄与距离成正比。

\textbf{7.} $u = 208\,276/10\,952 \approx \num{19.0}$~千米每百万年，
即 \qty{1.90}{cm} 每年。

\textbf{8.} 约 $2 \times 1.9 = \qty{3.8}{cm}$ 每年。

\textbf{9.} 是的：\qty{1.9}{cm} 每年略小于 \qty{2}{cm} 每
年。

\textbf{10.} $1270/19.0 \approx 67$ 百万年，与其
最老沉积物的年龄相同。

\textbf{11.} $191/9 \approx \num{21.2}$~千米每百万年，
\qty{2.1}{cm} 每年。

\textbf{12.} $1270/71 \approx \num{17.9}$~千米每百万年，
\qty{1.8}{cm} 每年。

\textbf{13.} $180\,329/9854 \approx \num{18.3}$~千米每百万年，
\qty{1.83}{cm} 每年。

\textbf{14.} $(1.90 - 1.83)/1.83 \approx \qty{4}{\percent}$：两个速率
在年龄和距离的不确定度范围内相符。

\textbf{15.} 一个年龄来自沉积物中的化石，另一个来自玄武岩的
磁性与磁倒转的时间标尺。一种方法中的误差
不会出现在另一种方法中：二者相符，
是对\omterm{def:g11:drift-to-spreading:spreading}{海底扩张}的有力检验。

\textbf{16.} $13.8 + 13.8 = \ang{27.6}$。

\textbf{17.} $111.2 \times \cos 30^\circ \approx \qty{96.3}{km}$ 每度；
$27.6 \times 96.3 \approx \qty{2660}{km}$。

\textbf{18.} $2660/19.0 \approx 140$ 百万年。

\textbf{19.} $2660/130 \approx \num{20.5}$~千米每百万年，约
\qty{2.0}{cm} 每年，略高于 \qty{1.9}{cm} 每年。最近 67 百万年
洋底以约 \qty{1.9}{cm} 每年的速率扩张，因此更早的张开略快：
20 号站点与边缘之间的洋底 $2660 - 1270 =
\qty{1390}{km}$，形成于 $130 - 67 = 63$ 百万年间，约
\qty{2.2}{cm} 每年。

\textbf{20.} \textbf{非洲沿岸南纬 \ang{30} 处最老的洋底
约 1.3 亿年}（按现今速率外推为 1.4 亿年）：
南大西洋于白垩纪在此处开始张开。
\end{solution}
